\section{Анализ предметной области и постановка задачи}
\thispagestyle{plain}

% TODO: ~14 страниц

\subsection{Современные системы интерактивного обучения}
% Quizlet, Kahoot, Mentimeter, российские: JoyTeka, OnlineTestPad, ЦОС МЭШ.
% Сравнительная таблица по: модель развёртывания, наличие real-time, геймификация,
% работа в РФ, открытый API/импорт, поддержка LMS.

\subsection{Геймификация в образовании}
% Подходы: Bartle taxonomy (1996), Octalysis Framework (Chou),
% Self-Determination Theory (Deci, Ryan).
% Эффективность: метаобзоры педагогических исследований по эффекту
% игровых механик на вовлечённость и удержание знаний.
% Интервальное повторение: алгоритмы Лейтнера, SuperMemo SM-2.

\subsection{Надёжность как предметная область и как требование к сервису}
% Двойственная роль:
% (а) Дисциплина <<Надёжность>>, преподаваемая на кафедре АСУ
%     (теория надёжности оборудования и программных систем,
%      нефтегазовый контекст: трубопроводы, насосное оборудование).
% (б) Требование к разрабатываемому сервису (отказоустойчивость,
%     резервирование, расчётно обоснованная доступность).

\subsection{Проблема работы образовательных сервисов в Российской Федерации}
% Ограничения иностранных облачных платформ (Firebase, AWS, Google Cloud)
% применительно к высшим учебным заведениям РФ.
% Альтернативы: серверная инфраструктура российских провайдеров
% (Selectel, Timeweb, Yandex Cloud, VK Cloud).
% Обоснование выбора <<собственный backend в РФ>> применительно к настоящей работе.

\subsection{Анализ научно-исследовательской работы и её ограничения}
% В рамках предшествующей научно-исследовательской работы (НИР)
% была разработана архитектура EduQuiz на стеке Flutter + Firebase.
% Выявленные ограничения:
% - зависимость от инфраструктуры Google;
% - минимальная геймификация (только лидерборд);
% - отсутствие интеграции с университетской LMS;
% - отсутствие экспериментальной верификации эффективности;
% - отсутствие расчёта показателей надёжности.

\subsection{Постановка задачи и требования}
% Функциональные требования: роли (преподаватель/студент/администратор),
% курсы, банки вопросов, интерактивные сессии, геймификация,
% импорт из edu.gubkin, аналитика.
% Нефункциональные требования: производительность, масштабируемость,
% надёжность, безопасность, удобство использования.
% Целевые показатели надёжности и производительности.

\textbf{Выводы по главе 1.}
% Резюме главы: что найдено, какой gap закрывает работа,
% что взято за основу.

\newpage
